\begin{abstract}

Semantic correspondence aims to establish pixel-level matches between semantically equivalent object parts across different images.
Despite extensive prior work, the task remains challenging due to large variations in viewpoint, scale, object deformation, and visual domain.
Recent advances in visual foundation models have shown that large-scale pretrained encoders learn rich semantic representations without explicit
correspondence supervision, suggesting strong potential for dense matching tasks.

In this work, we systematically study semantic correspondence using features extracted from pretrained visual foundation models.
We evaluate a training-free correspondence pipeline on the SPair-71k benchmark~\cite{min2019spair}, comparing three popular backbones: DINOv2~\cite{oquab2023dinov2}, DINOv3~\cite{simeoni2025dinov3}, and Segment Anything (SAM)~\cite{kirillov2023segment}.
We further analyze the effect of light fine-tuning of the final backbone layers using a contrastive objective, investigate an improved prediction
strategy based on windowed soft-argmax~\cite{zhang2024telling}, and assess the performance on another dataset, PF-Willow~\cite{ham2016}.

Our experiments show that DINO-based models significantly outperform SAM for semantic correspondence. In particular,
fine-tuned DINOv3 achieves a PCK@0.1 of \textbf{75.77\%}, slightly improving over fine-tuned DINOv2 (\textbf{74.89\%}), while fine-tuned DINOv2 attains the highest accuracy at
stricter thresholds, reaching \textbf{60.02\%} at PCK@0.05. Compared to training-free baselines, a small amount of fine-tuning yields significant improvements
across all thresholds. Additionally, windowed soft-argmax further enhances localization accuracy and robustness.

These results demonstrate that modern foundation models, when lightly adapted, provide a strong and scalable basis for semantic correspondence.
Project page: \url{https://github.com/Luffy65/Semantic-Correspondence}.

\end{abstract}